\documentclass[10pt,a4paper]{amsart}
\usepackage[margin=19mm]{geometry}
\usepackage{amsmath,amssymb,mathtools}
\usepackage[scr=rsfs,cal=euler]{mathalfa}
\usepackage{xfrac}
\usepackage[unicode,hidelinks,bookmarks=false]{hyperref}
\newcommand{\ie}{i.e.\ }
\newcommand{\cf}{cf.\ }
\newcommand{\resp}{respectively\ }
\newcommand{\Subsection}{\subsection}
\providecommand{\nobreakdash}{\nobreak\hbox{-}}
\setlength{\emergencystretch}{2em}
\multlinegap0pt
\usepackage{amssymb}
\usepackage{float}
\usepackage[shortlabels]{enumitem}
\newtheorem{thm}{Theorem}
\newtheorem{cor}{Corollary}
\newcommand{\F}{\mathbb{F}_q}
\title{Consecutive non-square non-primitive pairs in a finite field}
\author{Stephen D. Cohen}
\date{Source: arXiv:2604.21429v1 (CC BY 4.0)}
\begin{document}
\maketitle

\noindent\emph{Source and licence.} Consecutive non-square non-primitive pairs in a finite field. Version arXiv:2604.21429v1 (CC BY 4.0), distributed by arXiv under the Creative Commons Attribution 4.0 International licence, http://creativecommons.org/licenses/by/4.0/, as recorded in the arXiv licence metadata for this exact version and shown on the abstract page https://arxiv.org/abs/2604.21429v1. Full licence notice: \texttt{licenses/arxiv-2604.21429-CC-BY-4.0.txt}.\par\smallskip

\noindent\textit{Editorial scope.} Counted unit: the article's thm l7 (source0.tex lines 349--351). The article's complete original proof of it is delivered with it (source0.tex lines 353--368, one \texttt{proof} environment of 16 lines). No mathematics is written, repaired or re-derived: the statement and the proof are the article's own lines, in the article's own notation, and no retained line is substituted or re-flowed.  Retained uncounted as the source-local context the proof needs: lines 306--311; lines 315--317; lines 324--341.  Nothing was omitted from any retained span, so the package carries no omission record.  No retained line contains a citation, so the wrapper carries no bibliography and no bibliography entry.  No retained line contains a figure environment, an image-inclusion command or drawing code, so no image is delivered and the package declares no asset dependency.  Editorial formatting only, in addition: an AMS \texttt{amsart} preamble that restates the article's own declarations and macros used by the retained lines, namely the environments \texttt{thm}, \texttt{cor} on separate counters, together with the standard packages those lines need.  The retained environments print in this document as Theorem 1, Corollary 1 in order of appearance, whereas the article prints the same items with the numbering of its own full document.  The complete native source of the article as distributed in the arXiv e-print for this exact version is bundled unchanged as \texttt{sources/199023/source0.tex}.
\begin{thm}\label{thm:main}
Let $\ell$ be an odd prime dividing $q-1$. Then
\begin{equation}\label{mainbound}
N_\ell(q) \ge \frac{q - 4(\ell-1)^2 \sqrt{q} - (\ell-2)(1+\lambda(-1)) - 3}{4\ell^2}.
\end{equation}
\end{thm}

\begin{cor}\label{Npos}
Suppose $\ell$ is an odd prime divisor of $q-1$ such that $q > 16 \ell^4$. Then $\F$ contains an  NS$\ell$ pair.
\end{cor}

\begin{table}[h]\label{smalltable}
\centering 
\caption{Bounds for small $\ell$}
%\label{table}
\begin{tabular}{c|c|c|c}   
$\ell$ 
& min $q$ (by Th \ref{thm:main})
& min $q$ (by Cor \ref{Npos}) 
& min $q$ with $\theta_q \le 1/3$  \\[2pt]
 \hline
3  & 269          & 1,296        & 31 \\
5  & 4,117        & 10,000       & 771 \\
7  & 20,769       & 38,416       & 646,647 \\
11 & 160,045      & 234,256      & 70,673,696,523\\
13 & 331,829      & 456,976       & 28,215,721,625,272,767 \\
17 & 1,048,645    & 1,336,336     & 37,158,896,445,334,596,135,027\\
\end{tabular}
\end{table}

\begin{thm}\label{l7}
Let $q$ be an odd prime power such the least odd prime divisor of $q-1$ is $7$ and that $\theta_q \le 1/3$. Then $\F$  contains an  $NS7$ pair.
\end{thm}

\begin{proof}
As hypothesised, suppose  that $q$ is such that $\ell=7 \mid (q-1)$ but $3 \nmid (q-1)$ and $5 \nmid (q-1)$, and also that $\theta_q \le 1/3$.   Let $\omega_q$ denote the number of distinct prime divisors of $q-1$.
If $\omega_q \le 5$, then
\[
\theta_q \ge \frac{1}{2} \cdot \frac{6}{7} \cdot \frac{10}{11} \cdot \frac{12}{13} \cdot \frac{16}{17} > 0.338.
\]
Hence $\omega_q \ge 6$ and
\[
\theta_q \ge \frac{1}{2} \cdot \frac{6}{7} \cdot \frac{10}{11} \cdot \frac{12}{13} \cdot \frac{16}{17} \cdot \frac{18}{19} > 0.3206.
\]
It follows that
\[
q \ge 2 \cdot 7 \cdot 11 \cdot 13 \cdot 17 \cdot 19 + 1 = 646,647> 16\ell^4,
\]
and the result holds by Corollary  \ref{Npos} (as in the third row of Table \ref{smalltable}).
\end{proof}

\end{document}
